% Worked Examples: Concept 3.8.2 - Performance Specifications
\documentclass[11pt,a4paper]{article}
\usepackage{worked-examples-template}

\title{\textbf{Worked Examples}\\[0.5em]
\large Concept 3.8.2: Performance Specifications\\[0.3em]
\normalsize \coursesubtitle{} -- \coursetitle}
\author{Assoc.\ Prof.\ William Robertson \and Dr Sean McGowan}
\date{\today}

\begin{document}

\maketitle
\tableofcontents
\newpage

\section{Concept 3.8.2: Performance Specifications}

\subsection{Example 1: Computing Frequency Domain Performance Specifications}

Frequency domain specifications like maximum sensitivity, bandwidth, and peak frequency characterize system performance concisely.

\paragraph{Problem:} 
Consider a feedback system with process $P(s) = \frac{1}{(s+1)^3}$ and a PI controller $C(s) = k_p + \frac{k_i}{s}$ with gains $k_p = 0.6$ and $k_i = 0.5$.

(a) Calculate the loop transfer function $L(s)$ and the transfer function $G_{yr}(s) = \frac{PCF}{1+PC}$ from reference to output (assuming $F = 1$).

(b) Determine the complementary sensitivity function $T(s)$ and sensitivity function $S(s)$.

(c) Estimate the maximum sensitivity $M_s = \sup_\omega |S(\ii\omega)|$ and the frequency $\omega_{ms}$ at which it occurs.

(d) Estimate the bandwidth $\omega_b$ of $G_{yr}$, defined as the frequency where $|G_{yr}(\ii\omega_b)| = 1/\sqrt{2}$ (the -3 dB point).

(e) Relate the bandwidth to the system rise time using the approximation $T_r \omega_b \approx 3$.

\paragraph{Solution:}

\textbf{Step 1: Loop transfer function}

The PI controller is:
\[
C(s) = k_p + \frac{k_i}{s} = 0.6 + \frac{0.5}{s} = \frac{0.6s + 0.5}{s}
\]

The loop transfer function is:
\[
L(s) = P(s)C(s) = \frac{1}{(s+1)^3} \cdot \frac{0.6s + 0.5}{s} = \frac{0.6s + 0.5}{s(s+1)^3}
\]

\textbf{Step 2: Transfer function from reference to output}

With $F = 1$ (unity feedforward):
\[
G_{yr}(s) = \frac{PC}{1+PC} = \frac{L(s)}{1+L(s)} = T(s)
\]

The complementary sensitivity function is:
\[
T(s) = \frac{0.6s + 0.5}{s(s+1)^3 + 0.6s + 0.5}
\]

Expanding the denominator:
\[
s(s+1)^3 = s(s^3 + 3s^2 + 3s + 1) = s^4 + 3s^3 + 3s^2 + s
\]
\[
s(s+1)^3 + 0.6s + 0.5 = s^4 + 3s^3 + 3s^2 + 1.6s + 0.5
\]

Therefore:
\[
T(s) = G_{yr}(s) = \frac{0.6s + 0.5}{s^4 + 3s^3 + 3s^2 + 1.6s + 0.5}
\]

\textbf{Step 3: Sensitivity function}

Using $S = 1 - T$:
\[
S(s) = \frac{s^4 + 3s^3 + 3s^2 + s}{s^4 + 3s^3 + 3s^2 + 1.6s + 0.5}
\]

Alternatively, directly:
\[
S(s) = \frac{1}{1+L(s)} = \frac{s(s+1)^3}{s^4 + 3s^3 + 3s^2 + 1.6s + 0.5}
\]

\textbf{Step 4: Maximum sensitivity $M_s$}

The maximum sensitivity occurs at a frequency where $|S(\ii\omega)|$ peaks. For this system, we need to evaluate:
\[
|S(\ii\omega)| = \left|\frac{\ii\omega(\ii\omega + 1)^3}{(\ii\omega)^4 + 3(\ii\omega)^3 + 3(\ii\omega)^2 + 1.6\ii\omega + 0.5}\right|
\]

At $\omega = 0$: $|S(0)| = \frac{0}{0.5} = 0$

At $\omega \to \infty$: $|S(\ii\omega)| \to \frac{\omega \cdot \omega^3}{\omega^4} = 1$

The peak typically occurs at a mid-range frequency. By numerical evaluation or using control system software:

At $\omega = 1$ rad/s:
\[
S(\ii) = \frac{\ii(1+\ii)^3}{(\ii)^4 + 3(\ii)^3 + 3(\ii)^2 + 1.6\ii + 0.5}
\]

Computing $(1+\ii)^3$:
\[
(1+\ii)^2 = 1 + 2\ii - 1 = 2\ii
\]
\[
(1+\ii)^3 = (1+\ii) \cdot 2\ii = 2\ii + 2\ii^2 = 2\ii - 2 = -2 + 2\ii
\]

So:
\[
\ii(1+\ii)^3 = \ii(-2 + 2\ii) = -2\ii - 2 = -2 - 2\ii
\]

The denominator at $s = \ii$:
\[
(\ii)^4 + 3(\ii)^3 + 3(\ii)^2 + 1.6\ii + 0.5 = 1 - 3\ii - 3 + 1.6\ii + 0.5 = -1.5 - 1.4\ii
\]

Therefore:
\[
S(\ii) = \frac{-2 - 2\ii}{-1.5 - 1.4\ii}
\]

Computing the magnitude:
\[
|S(\ii)| = \frac{\sqrt{4 + 4}}{\sqrt{2.25 + 1.96}} = \frac{\sqrt{8}}{\sqrt{4.21}} = \frac{2.828}{2.052} = 1.378
\]

Testing $\omega = 0.8$ rad/s would give a value closer to the peak. Through numerical optimization, we find:
\[
M_s \approx 1.4, \quad \omega_{ms} \approx 0.9 \text{ rad/s}
\]

\textbf{Step 5: Bandwidth calculation}

The bandwidth $\omega_b$ is where $|G_{yr}(\ii\omega_b)| = 1/\sqrt{2} \approx 0.707$.

At DC: $|G_{yr}(0)| = \frac{0.5}{0.5} = 1$

As frequency increases, the magnitude decreases. By numerical evaluation or graphical methods:
\[
\omega_b \approx 1.0 \text{ rad/s}
\]

\textbf{Step 6: Rise time estimation}

Using the rise time-bandwidth product approximation:
\[
T_r \approx \frac{3}{\omega_b} = \frac{3}{1.0} = 3 \text{ seconds}
\]

This can be verified by simulating the step response and measuring the time from 10\% to 90\% of the final value.

\textbf{Key insights:}
\begin{itemize}
\item $M_s \approx 1.4$ indicates reasonable robustness (typically want $M_s < 2$ for good margins)
\item The stability margin is approximately $s_m = 1/M_s \approx 0.71$, meaning the Nyquist curve stays at least 0.71 units away from the critical point $-1$
\item Bandwidth $\omega_b \approx 1.0$ rad/s indicates the frequency range where the system responds effectively to reference commands
\item The rise time-bandwidth product provides a quick estimate: faster responses require larger bandwidth
\item The peak in $|S|$ near $\omega_{ms}$ indicates the frequency range where disturbances are most amplified (though still less than 1.4×)
\end{itemize}

\subsection{Example 2: Error Coefficients and Steady-State Tracking Performance}

Error coefficients characterize how well a system tracks different types of reference signals at steady state.

\paragraph{Problem:}
Consider a unity feedback system with process $P(s) = \frac{1}{s(s+1)}$ (a type-1 system due to the pole at $s=0$). Two controllers are being considered:
\begin{itemize}
\item Controller A: $C_A(s) = 5$ (proportional only)
\item Controller B: $C_B(s) = 2(1 + \frac{1}{s})$ (PI controller)
\end{itemize}

(a) For each controller, derive the error transfer function $G_{er}(s) = \frac{1}{1+PC}$ from reference $r$ to tracking error $e = r - y$.

(b) Expand $G_{er}(s)$ as a power series in $s$ to find the error coefficients: $G_{er}(s) = e_0 + e_1 s + e_2 s^2 + \cdots$

(c) Determine the steady-state error for:
\begin{itemize}
\item Step input: $r(t) = R_0$ (so $\dot{r} = 0$, $\ddot{r} = 0$)
\item Ramp input: $r(t) = v_0 t$ (so $\dot{r} = v_0$, $\ddot{r} = 0$)
\item Parabolic input: $r(t) = a_0 t^2$ (so $\ddot{r} = 2a_0$)
\end{itemize}

(d) Compare the tracking performance of the two controllers.

\paragraph{Solution:}

\textbf{Step 1: Error transfer function for Controller A}

Loop transfer function:
\[
L_A(s) = P(s)C_A(s) = \frac{1}{s(s+1)} \cdot 5 = \frac{5}{s(s+1)}
\]

Error transfer function:
\[
G_{er,A}(s) = \frac{1}{1+L_A(s)} = \frac{1}{1 + \frac{5}{s(s+1)}} = \frac{s(s+1)}{s(s+1) + 5} = \frac{s^2 + s}{s^2 + s + 5}
\]

\textbf{Step 2: Power series expansion for Controller A}

To find the error coefficients, we expand $G_{er,A}(s)$ for small $s$:
\[
G_{er,A}(s) = \frac{s^2 + s}{s^2 + s + 5} = \frac{s^2 + s}{5} \cdot \frac{1}{1 + \frac{s^2 + s}{5}}
\]

Using the expansion $\frac{1}{1+x} \approx 1 - x + x^2 - \cdots$ for small $x$:
\[
G_{er,A}(s) \approx \frac{s^2 + s}{5}\left(1 - \frac{s^2 + s}{5} + \cdots\right)
\]
\[
\approx \frac{s + s^2}{5} - \frac{(s + s^2)^2}{25} + \cdots
\]
\[
\approx \frac{s}{5} + \frac{s^2}{5} + O(s^2)
\]

Therefore:
\[
e_0 = 0, \quad e_1 = \frac{1}{5}, \quad e_2 = \frac{1}{5}
\]

\textbf{Step 3: Error transfer function for Controller B}

Loop transfer function:
\[
L_B(s) = P(s)C_B(s) = \frac{1}{s(s+1)} \cdot \frac{2(s+1)}{s} = \frac{2(s+1)}{s^2(s+1)} = \frac{2}{s^2}
\]

Error transfer function:
\[
G_{er,B}(s) = \frac{1}{1 + \frac{2}{s^2}} = \frac{s^2}{s^2 + 2}
\]

\textbf{Step 4: Power series expansion for Controller B}

For small $s$:
\[
G_{er,B}(s) = \frac{s^2}{s^2 + 2} = \frac{s^2}{2} \cdot \frac{1}{1 + \frac{s^2}{2}}
\]

Using $\frac{1}{1+x} \approx 1 - x + \cdots$:
\[
G_{er,B}(s) \approx \frac{s^2}{2}\left(1 - \frac{s^2}{2} + \cdots\right) \approx \frac{s^2}{2} + O(s^4)
\]

Therefore:
\[
e_0 = 0, \quad e_1 = 0, \quad e_2 = \frac{1}{2}
\]

\textbf{Step 5: Steady-state errors}

The steady-state error for a reference signal $r(t)$ is:
\[
e_{ss} = e_0 r(t) + e_1 \dot{r}(t) + e_2 \ddot{r}(t) + \cdots
\]

\textbf{For a step input $r(t) = R_0$:}

Controller A: $e_{ss,A} = 0 \cdot R_0 + \frac{1}{5} \cdot 0 + \cdots = 0$

Controller B: $e_{ss,B} = 0 \cdot R_0 + 0 \cdot 0 + \cdots = 0$

Both controllers have zero steady-state error for step inputs (both are type-1 systems).

\textbf{For a ramp input $r(t) = v_0 t$:}

Controller A: $e_{ss,A} = 0 + \frac{1}{5} \cdot v_0 + 0 = \frac{v_0}{5}$

Controller B: $e_{ss,B} = 0 + 0 \cdot v_0 + 0 = 0$

Controller B (PI) has zero steady-state error for ramps, while Controller A has a constant lag error.

\textbf{For a parabolic input $r(t) = a_0 t^2$:}

Controller A: $e_{ss,A} = 0 + \frac{1}{5} \cdot 2a_0 t + \frac{1}{5} \cdot 2a_0 = \frac{2a_0 t}{5} + \frac{2a_0}{5}$

The error grows linearly with time (unbounded).

Controller B: $e_{ss,B} = 0 + 0 + \frac{1}{2} \cdot 2a_0 = a_0$

Controller B has a constant error (bounded but non-zero).

\textbf{Step 6: Comparison and interpretation}

\begin{center}
\begin{tabular}{|l|c|c|}
\hline
\textbf{Input} & \textbf{Controller A (P)} & \textbf{Controller B (PI)} \\
\hline
Step ($R_0$) & 0 & 0 \\
Ramp ($v_0 t$) & $v_0/5$ & 0 \\
Parabolic ($a_0 t^2$) & $\infty$ (grows) & $a_0$ (constant) \\
\hline
\end{tabular}
\end{center}

\textbf{Key insights:}
\begin{itemize}
\item The system type (number of integrators in the loop) determines which polynomial inputs can be tracked with zero steady-state error
\item The process $P(s) = 1/[s(s+1)]$ already has one integrator, making it type-1 without any controller
\item Adding PI control (Controller B) introduces a second integrator, making the loop type-2
\item Type-1 systems track steps with zero error and ramps with constant error
\item Type-2 systems track both steps and ramps with zero error, and parabolas with constant error
\item Higher-order terms ($e_2, e_3, \ldots$) determine the error for higher-order polynomial inputs
\item Error coefficients provide insight without needing to simulate each input type
\item Low-frequency behavior (small $s$) corresponds to long-time steady-state behavior
\end{itemize}
\end{document}
